\MathBookSourcePage{249}
\subsection{习题}
\label{sec:ch08-exercises}
\setcounter{exerciseitem}{0}

\begin{MBExercise}
\label{exr:ch08-01}
证明对 $2<p<\infty$, 定理~\ref{thm:ch08-main-max-norm-stability} 中可将 $\infty$ 换成 $p$. 提示: 使用 Hölder 不等式, 使式~\eqref{eq:ch08-green-error-action-bound} 中保留 $\nabla u$ 的一个加权 $L^p$ 范数; 对式~\eqref{eq:ch08-localization-identity} 的第 $p$ 次幂积分, 再应用 Fubini(参见 Rannacher and Scott 1982).
\end{MBExercise}

\begin{MBExercise}
\label{exr:ch08-02}
证明推论~\ref{cor:ch08-main-max-norm-error}. 提示: 写成 $u-u_h=u-I^hu+(I^hu-u_h)$, 并将定理~\ref{thm:ch08-main-max-norm-stability} 用于以 $I^hu-u$ 代替 $u$ 的情形.
\end{MBExercise}

\begin{MBExercise}
\label{exr:ch08-03}
证明式~\eqref{eq:ch08-weight-element-comparability}. 常数如何依赖于 $\kappa$ 和 $h$?
\end{MBExercise}

\begin{MBExercise}
\label{exr:ch08-04}
证明式~\eqref{eq:ch08-weight-infinity-bound}. 常数如何依赖于 $\kappa$ 和 $h$?
\end{MBExercise}

\begin{MBExercise}
\label{exr:ch08-05}
证明式~\eqref{eq:ch08-weight-derivative-bound}. 常数如何依赖于 $\kappa$ 和 $h$? 提示: 先考虑 $\lambda=1$ 的情形, 再用归纳法.
\end{MBExercise}

\begin{MBExercise}
\label{exr:ch08-06}
证明式~\eqref{eq:ch08-weighted-interpolation-estimates}. 常数如何依赖于 $\kappa$ 和 $h$? 提示: 使用式~\eqref{eq:ch08-weight-element-comparability}.
\end{MBExercise}

\begin{MBExercise}
\label{exr:ch08-07}
证明式~\eqref{eq:ch08-smoothed-delta-properties}. 常数如何依赖于 $k$? 提示: 参见 Scott (1976).
\end{MBExercise}

\begin{MBExercise}
\label{exr:ch08-08}
证明式~\eqref{eq:ch08-weighted-green-second-derivative} 由引理~\ref{lem:ch08-weighted-primal-regularity} 推出. 提示: 使用式~\eqref{eq:ch08-weight-infinity-bound}.
\end{MBExercise}

\begin{MBExercise}
\label{exr:ch08-09}
详细推导式~\eqref{eq:ch08-weighted-primal-energy-bound}, 并结合 Schwarz 不等式和 Young 不等式~\eqref{eq:ch00-young-inequality-delta}, 适当拆分权函数.
\end{MBExercise}

\begin{MBExercise}
\label{exr:ch08-10}
作注~\ref{rem:ch08-index-engineering} 所建议的变量替换, 验证其中不等式两端 $L$ 的幂次相同.
\end{MBExercise}

\begin{MBExercise}
\label{exr:ch08-11}
证明二维中的 Sobolev 嵌入 $W_1^1(\Omega)\subset L^2(\Omega)$. 提示: 写成
\[
\MathBookFormulaID{formula-ch08-exercise-11-sobolev-representation}
u(x_1,x_2)^2
=\int_{-\infty}^{x_1}\frac{\partial u}{\partial x_1}(t,x_2)\,dt
\int_{-\infty}^{x_2}\frac{\partial u}{\partial x_2}(x_1,s)\,ds,
\]
对 $x$ 积分并使用 Fubini.
\end{MBExercise}

\begin{MBExercise}
\label{exr:ch08-12}
对 $u\in\mathring W_q^1(\Omega)$, 证明 Sobolev 不等式
\[
\MathBookFormulaID{formula-ch08-exercise-12-sobolev-inequality}
\|u\|_{L^p(\Omega)}\leq c_p\|u\|_{W_q^1(\Omega)},
\qquad 1+\frac2p=\frac2q,
\]
其中为二维情形. 当 $p\to\infty$ 时 $c_p$ 如何变化? 提示: 写成
\[
\MathBookFormulaID{formula-ch08-exercise-12-power-identity}
\|u\|_{L^p(\Omega)}^p=\|u^{p/2}\|_{L^2(\Omega)}^2
\leq C\|u^{p/2}\|_{W_1^1(\Omega)}^2,
\]
使用习题~\ref{exr:ch08-11}, 展开 $\nabla(u^{p/2})$, 并应用 Hölder 不等式.
\end{MBExercise}

\begin{MBExercise}
\label{exr:ch08-13}
在定理~\ref{thm:ch08-main-max-norm-stability} 的条件下证明式~\eqref{eq:ch08-lp-lower-order-error}. 提示: 以 $f=|u-u_h|^{p-1}\operatorname{sign}(u-u_h)$ 解式~\eqref{eq:ch08-dual-lp-problem}, 再沿用后续估计.
\end{MBExercise}

\begin{MBExercise}
\label{exr:ch08-14}
式~\eqref{eq:ch08-dual-lp-regularity} 的证明中 $r$ 应如何选择? 提示: 令 $q$ 和 $r$ 由 Sobolev 不等式~\ref{exr:ch04-11} 联系.
\end{MBExercise}

\begin{MBExercise}
\label{exr:ch08-15}
对二维和三维拟一致剖分上的 Lagrange、Hermite 和 Argyris 单元, 以及拟一致矩形剖分上的张量积与 Serendipity 单元, 证明式~\eqref{eq:ch08-weighted-interpolation-estimates}. 提示: 参见定理~\ref{thm:ch04-local-tensor-product-interpolation} 的证明并使用式~\eqref{eq:ch08-weight-element-comparability}.
\end{MBExercise}

\begin{MBExercise}
\label{exr:ch08-16}
对第~\ref{chap:ch03-finite-element-space-construction} 章在拟一致剖分上介绍的所有分片多项式有限元, 证明式~\eqref{eq:ch08-weighted-inverse-estimate}. 提示: 参见引理~\ref{lem:ch04-local-inverse-estimate} 的证明并使用式~\eqref{eq:ch08-weight-element-comparability}.
\end{MBExercise}

\MathBookSourcePage{251}
\begin{MBExercise}
\label{exr:ch08-17}
设 $A=(a_{ij})$ 为 $d\times d$ 对称矩阵, 满足
\[
\MathBookFormulaID{formula-ch08-exercise-17-spectral-bounds}
C_a|\xi|^2\leq\sum_{i,j=1}^{d}a_{ij}\xi_i\xi_j\leq M|\xi|^2
\qquad\forall\xi\in\mathbb R^d,
\]
其中 $C_a>0$. 证明若 $AX=\lambda X$ 对某个 $X\neq0$ 成立, 则 $C_a\leq\lambda\leq M$. 再证明 $C_a$(相应地 $M$)可取为特征值的最小值(相应地最大值). 最后证明
\[
\MathBookFormulaID{formula-ch08-exercise-17-bilinear-bound}
\sum_{i,j=1}^{d}a_{ij}\xi_i\nu_j\leq M|\xi||\nu|
\qquad\forall\xi,\nu\in\mathbb R^d.
\]
\end{MBExercise}

\begin{MBExercise}
\label{exr:ch08-18}
对伴随问题~\eqref{eq:ch08-adjoint-problem} 证明定理~\ref{thm:ch08-main-max-norm-stability}, 其中 $u_h$ 由 $a(v,u_h)=(f,v)$ 对所有 $v\in V_h$ 定义. 提示: 考察
\[
\MathBookFormulaID{formula-ch08-exercise-18-adjoint-form}
a(u,v):=\int_\Omega\left\{\sum_{i,j=1}^{d}a_{ij}
\frac{\partial u}{\partial x_i}\frac{\partial v}{\partial x_j}
+\sum_{i=1}^{d}b_i\frac{\partial v}{\partial x_i}u+b_0uv\right\}\,dx,
\]
并对论证作适当修改.
\end{MBExercise}

\begin{MBExercise}
\label{exr:ch08-19}
设 $G_h^z\in V_h$ 为二维区域 $\Omega$ 上某个 $m$ 次分片多项式空间 $V_h$ 中由
\[
\MathBookFormulaID{formula-ch08-exercise-19-discrete-green}
a(G_h^z,v)=v(z)
\qquad\forall v\in V_h
\]
定义的离散 Green 函数. 证明
\[
\MathBookFormulaID{formula-ch08-exercise-19-discrete-green-bound}
G_h^z(z)\leq C|\log h|.
\]
提示: 因为 $G_h^z\in V_h$, 有 $G_h^z(z)=a(G_h^z,G_h^z)$. 应用离散 Sobolev 不等式~\ref{lem:ch04-discrete-sobolev-inequality} 证明
\[
\MathBookFormulaID{formula-ch08-exercise-19-energy-bound}
G_h^z(z)\leq C|\log h|^{1/2}a(G_h^z,G_h^z)^{1/2}.
\]
参见 Bramble (1966).
\end{MBExercise}

\begin{MBExercise}
\label{exr:ch08-20}
设 $G_h^z\in V_h$ 为二维区域 $\Omega$ 上某个满足 Dirichlet 边界条件的 $m$ 次分片多项式空间 $V_h$ 中由
\[
\MathBookFormulaID{formula-ch08-exercise-20-discrete-green}
a(G_h^z,v)=v(z)
\qquad\forall v\in V_h
\]
定义的离散 Green 函数. 再假设 $z$ 到 $\partial\Omega$ 的距离为 $O(h)$. 证明
\[
\MathBookFormulaID{formula-ch08-exercise-20-discrete-green-uniform-bound}
G_h^z(z)\leq C.
\]
提示: 参见习题~\ref{exr:ch08-19}. 因为 $G_h^z=0$ 于 $\partial\Omega$, 有
\[
\MathBookFormulaID{formula-ch08-exercise-20-boundary-inverse-bound}
G_h^z(z)\leq Ch|G_h^z|_{W_\infty^1(\Omega)}.
\]
改用逆不等式而非离散 Sobolev 不等式~\ref{lem:ch04-discrete-sobolev-inequality} 完成证明. 参见 Dragancescu, Dupont and Scott (2002).
\end{MBExercise}
